\begin{frame}[allowframebreaks]{유도 템플릿 ②: SVM KKT --- ``C가 실제로 바꾸는 것''}
  Week 5.2 소프트 마진 SVM에서 ``C가 커지면 마진이 커진다''는
  \textbf{잘못된 직감}을 시험에서 확인하는 문제가 자주 나온다.
  \vskip0.4em
  \textbf{문제} --- 1차원 4개 점:
  \small
  \begin{tabular}{@{}clcl@{}}
    \toprule
    점 & $x$ & $y$ & 역할 \\
    \midrule
    A & $-1$ & $+1$ & \\
    B & $-1$ & $+1$ & A와 같은 위치(다중점) \\
    C & $+2$ & $-1$ & \\
    E & $-0.2$ & $-1$ & \textbf{이상치} --- $-1$인데 $+1$ 쪽에
      붙어 있음 \\
    \bottomrule
  \end{tabular}
  \vskip0.3em
  목적함수
  \[ \min_{w,b}\ \tfrac12 w^2 + C\sum_i \xi_i
     \qquad (\text{약분 } y_i(wx_i+b) \ge 1-\xi_i,\ \xi_i \ge 0) \]
  A·B가 같은 $(x,y)$이므로 대칭으로 $\alpha_A=\alpha_B\triangleq a$.
  \vskip0.4em
  \kb{KKT 세 조건} + 정합성:
  \begin{enumerate}
    \item \kb{$0 \le \alpha_i \le C$} --- $\alpha$는 ``점 $i$에게
      쓰인 위반 예산''. C가 전체 예산이면 $\alpha$는 각 점에
      배정된 몫.
    \item \kb{$\alpha_i \xi_i = 0$} --- 마진 위반($\xi_i>0$)을
      하는 점은 $\alpha$가 상한 C에 \textbf{꽉 차 있다}.
    \item \kb{$\alpha_i(y_if(x_i)-1+\xi_i)=0$}
      \begin{itemize}
        \item $\alpha_i=0$: 마진 바깥($y_if(x_i)\ge 1$), 해에
          영향 없음.
        \item $0<\alpha_i<C$: \textbf{마진 위에 정확히}
          ($y_if(x_i)=1$).
        \item $\alpha_i=C$: 마진 안 또는 오분류
          ($y_if(x_i)\le 1$).
      \end{itemize}
  \end{enumerate}
  \vskip0.3em
  \begin{block}{정합성 (stationarity)}
    \[ w = \sum_i \alpha_i y_i x_i, \qquad \sum_i \alpha_i y_i = 0 \]
    이 문제에선 클래스 2:2라 $\alpha_A+\alpha_B = \alpha_C+
    \alpha_E$. 이제 $C$를 0에서부터 올리며 각 구간의 해를 구한다.
    \kb{이 4구간 풀이(구간 경계 = ``어떤 $\alpha$가 상한/0에
    닿는가'')가 템플릿}.
  \end{block}
  \vskip0.3em
  \scriptsize
  \begin{tabular}{@{}cllll@{}}
    \toprule
    구간 & $C$ 범위 & $\alpha$ & $w, b$ & $\|w\|^2$, 마진
      $2/\|w\|$ \\
    \midrule
    I & $C \le \tfrac{10}{57}\approx0.175$
      & 전부 $=C$ & $(-3.8C,\ 1.9C)$ & $14.44\,C^2$; 마진
      $\tfrac{2}{3.8C}$ (52.6$\to$3.0) \\
    II & $\tfrac{10}{57}<C<\tfrac56$
      & $\alpha_E=C$, $a=\tfrac19+\tfrac{11}{30}C$,
        $\alpha_C=\tfrac29-\tfrac4{15}C$
      & \textbf{$(-\tfrac23,\ \tfrac13)$ 고정} &
      \textbf{4/9 고정}; 마진 \textbf{3.0 고정} \\
    III & $\tfrac56<C<\tfrac{25}{8}$
      & $\alpha_E=C$, $a=C/2$, $\alpha_C=0$ & $(-0.8C,\
        1-0.8C)$ & $0.64\,C^2$; 마진 $\tfrac{2.5}{C}$
      (3.0$\to$0.8) \\
    IV & $C \ge \tfrac{25}{8}=3.125$
      & $a=\tfrac{25}{16}$, $\alpha_E=\tfrac{25}{8}$,
        $\alpha_C=0$
      & \textbf{$(-\tfrac52,\ -\tfrac32)$ 고정} &
      \textbf{6.25 고정}; 마진 \textbf{0.8 고정} \\
    \bottomrule
  \end{tabular}
  \vskip0.4em
  {\footnotesize 구간 경계: I$\to$II는 A·B가 마진 위에 올 때까지
  $f(-1)=5.7C=1 \Rightarrow C=10/57$. II$\to$III은 $\alpha_C=0$이
  되는 $C=5/6$. III$\to$IV은 E가 마진 위에 오기까지
  $y_Ef(x_E)=0.64C-1=1 \Rightarrow C=25/8$. 참고: E는
  $C<25/16=1.5625$에서는 아직 \textbf{오분류} 상태.}
\end{frame}
